\documentclass[10pt,a4paper]{amsart}
\usepackage[margin=19mm]{geometry}
\usepackage{amsmath,amssymb,mathtools}
\usepackage[scr=rsfs,cal=euler]{mathalfa}
\usepackage{xfrac}
\usepackage[unicode,hidelinks,bookmarks=false]{hyperref}
\newcommand{\ie}{i.e.\ }
\newcommand{\cf}{cf.\ }
\newcommand{\resp}{respectively\ }
\newcommand{\Subsection}{\subsection}
\providecommand{\nobreakdash}{\nobreak\hbox{-}}
\setlength{\emergencystretch}{2em}
\multlinegap0pt
\usepackage{bm}
\usepackage{bbm}
\usepackage{dsfont}
\usepackage{xspace}
\usepackage{cancel}
\usepackage{mathtools}
\usepackage{amsthm}
\newtheorem{theorem}{Theorem}
\title[Total Variation Distance Estimation through Domain Reduction]{Total Variation Distance Estimation through Domain Reduction (Theorem 5)}
\author{Arnab Bhattacharyya, Graham Cormode, Yucheng Fu, Kuldeep S. Meel}
\date{Source: arXiv:2609.18707v2 (CC BY 4.0)}
\begin{document}
\maketitle

\noindent\emph{Source and licence.} Total Variation Distance Estimation through Domain Reduction. Version arXiv:2609.18707v2 (CC BY 4.0), distributed under the Creative Commons Attribution 4.0 International licence, as recorded in the arXiv licence metadata for this exact version and shown on the abstract page https://arxiv.org/abs/2609.18707v2. Full rights notice: licenses/arxiv-2609.18707-CC-BY-4.0.txt.\par\smallskip

\noindent\textit{Editorial scope.} Counted unit: the Theorem printed as Theorem 5 of the article (source0.tex lines 997--1004, source label \texttt{thm:wta\_fpras}), delivered together with the article's own complete original proof of it (source0.tex lines 1006--1036, one \texttt{proof} environment of 31 lines). No mathematics is written, repaired or re-derived: statement and proof are the article's own text line for line. Required context retained uncounted: thm:pc\_fpras (lines 901--914). Every definition, display and label of those lines is retained unchanged. Omitted from this package and not redistributed: 1--900, 915--996, 1005--1005, 1037--1050. Nothing inside a retained excerpt is omitted or altered: every retained body is delivered exactly as the article's lines are written, so no omission receipt of any kind accompanies this package. Citations: the retained excerpts carry no citation command at all, so no bibliography entry of the article is transcribed and no reference is redistributed. Reference adaptation: none was needed and none was made; every reference inside the delivered unit resolves inside it, so no unresolved reference or citation is produced. Printed slips kept literal: no printed word, symbol, hypothesis or number of the counted statement or of its proof was altered. Layout and class: the article's own class is \texttt{article} with packages including nameref, hyperref, amsmath, amssymb, amsthm, geometry, algorithm, algorithmic, tikz; the wrapper is the lane's pinned \texttt{amsart} editorial wrapper at 10pt on A4, which restates the theorem-like environments the retained text needs and adds the article's own macro definitions that the retained text needs (none), with extra wrapper packages (bm, bbm, dsfont, xspace, cancel, mathtools). The delivered numbering is the unit's own. Assets: no retained span contains a figure environment, a graphics-inclusion command, a PDF-page inclusion, a table or a TikZ picture; no image file is copied into this package, the package declares no asset dependency and no image file is redistributed with it. Licence: version arXiv:2609.18707v2 of this article is distributed under the Creative Commons Attribution 4.0 International licence, as recorded in the arXiv licence metadata for this exact version and shown on the abstract page https://arxiv.org/abs/2609.18707v2. A full rights notice is delivered at licenses/arxiv-2609.18707-CC-BY-4.0.txt. Characters: every retained excerpt is pure ASCII, so the delivered engine, XeLaTeX with its default Latin Modern text font, reports no missing character.
\begin{theorem}[Probabilistic Circuit TV Approximation]
\label{thm:pc_fpras}
Given $\epsilon, \eta \in (0, 1)$ and smooth, decomposable probabilistic circuits $C_P$ and $C_Q$ computing $P$ and $Q$, respectively, and structured with respect to a common v-tree, let $W$, $q$, and $L\geq1$ be as defined above. The bottom-up coreset propagation outputs the estimator $\widetilde{D}$ defined above, satisfying $(1-\epsilon)d_{\mathrm{TV}}(P,Q) \le \widetilde{D} \le (1+\epsilon)d_{\mathrm{TV}}(P,Q)$ with probability at least $1-\eta$. Let $M$ denote the maximum number of rows retained in any
internal-region coreset, and
let $|C|:=|C_P|+|C_Q|$, where each circuit size counts its gates and wires. The arithmetic running time is
\[
\widetilde{O}\!\left((q+M)|C|+L\left[W(q+M)^2+W^\omega\right]\right),
\quad\text{where }
M=O\!\left(
\frac{WL^2\log(2W)}{\epsilon^2}
\log\frac{L}{\eta}
\right).
\]
\end{theorem}

\begin{theorem}[FPRAS for TV-distance between weighted tree automata]\label{thm:wta_fpras}
Given $\varepsilon,\eta\in(0,1)$, two distributions $P$ and $Q$ induced by two nonnegative weighted tree automata on the same tree $\mathcal{T}$, there exists an FPRAS that outputs a real number $\widehat d$ such that
$(1-\varepsilon)d_{\mathrm{TV}}(P,Q)\leq\widehat d\leq(1+\varepsilon)d_{\mathrm{TV}}(P,Q)$, with probability at least $1-\eta$ in time
$$
\widetilde{O}\!\left(|V|d^2\left(q+d+\frac{d^2|V|^2}{\varepsilon^2}\log\frac{|V|}{\eta}\right)^2+|V|d^{2\omega}\right),
$$
where $q:=\max_{i\in[|L|]}|\Omega_i|$ and $d:=\max_{u\in V} d_u$ are the maximum leaf-domain size and joint node-interface dimension, respectively.
\end{theorem}

\begin{proof}
At every node $u$, we pad each model's message vector with zero coordinates to the common global dimension $d$. The corresponding leaf tables and transition tensors are padded with zeros. This padding does not change either induced distribution. For notational simplicity, we henceforth reuse $h_u^R$ and $T_u^R$ to denote the corresponding zero-padded messages and tensors. Since $Z_R$ can be computed exactly by the preceding bottom-up dynamic program, we append a unary root sum gate with weight $Z_R^{-1}$, whose output is $R(x)=Z_R^{-1}f_R(x)$. In the expanded root feature space, the two-dimensional query $a_{\mathrm{TV}}$ defined above is extended by zero coefficients on all other root-scope gate coordinates.

We next represent every bilinear gate $\mathcal B_u^R$ exactly as a layer of scalar product gates followed by a layer of scalar sum gates. For every node $u\in V$ and every $i\in[d]$, let $g_{u,i}$ denote the gate corresponding to the $i$-th coordinate of the padded message at node $u$. At a leaf $\ell$, $g_{\ell,i}$ is a leaf gate satisfying
$$
R_{g_{\ell,i}}(x_{S_\ell}) = \bigl[h_\ell^R(x_{S_\ell})\bigr]_i
$$
under each parameterization $R\in\{P,Q\}$.

Now let $u$ be an internal node. Recall that its bilinear gate is given coordinatewise by
$$
\bigl[\mathcal B_u^R(p,q)\bigr]_i = \sum_{j=1}^{d}\sum_{k=1}^{d} T_{u,i,j,k}^R\,p_jq_k, \qquad i\in[d].
$$
For every $j\in[d]$ and $k\in[d]$, introduce a product gate $p_{u,j,k}$ whose children are $g_{u_{\mathrm L},j}$ and $g_{u_{\mathrm R},k}$. Under parameterization $R$, it computes
$$
R_{p_{u,j,k}}(x_{S_u}) = R_{g_{u_{\mathrm L},j}}(x_{S_{u_{\mathrm L}}}) R_{g_{u_{\mathrm R},k}}(x_{S_{u_{\mathrm R}}}).
$$

For every $i\in[d]$, define $g_{u,i}$ to be a sum gate whose children are all product gates $p_{u,j,k}$, with edge weight $T_{u,i,j,k}^R$ under parameterization $R$. Thus,
$$
R_{g_{u,i}}(x_{S_u}) = \sum_{j=1}^{d}\sum_{k=1}^{d} T_{u,i,j,k}^R\, R_{p_{u,j,k}}(x_{S_u}).
$$
Hence this two-layer subcircuit computes the bilinear gate:
$$
\begin{bmatrix} R_{g_{u,1}}(x_{S_u})\\ \vdots\\ R_{g_{u,d}}(x_{S_u}) \end{bmatrix} = \mathcal B_u^R\!\left(h_{u_{\mathrm L}}^R(x_{S_{u_{\mathrm L}}}), h_{u_{\mathrm R}}^R(x_{S_{u_{\mathrm R}}})\right) = h_u^R(x_{S_u}).
$$

Each product gate $p_{u,j,k}$ is decomposable because its two children have the disjoint scopes $S_{u_{\mathrm L}}$ and $S_{u_{\mathrm R}}$. Each sum gate $g_{u,i}$ is smooth because all of its children have scope $S_u=S_{u_{\mathrm L}}\sqcup S_{u_{\mathrm R}}$. All product gates respect $\mathcal T$, so the two expanded circuits are structured with respect to the same v-tree. Therefore, Theorem~\ref{thm:pc_fpras} gives the stated approximation guarantee.

For the running time, let $I:=|V\setminus L|$ be the number of internal nodes. If $I=0$, the TV-distance can be computed exactly by enumerating the single leaf domain, so assume $I\geq1$. Each expanded circuit has at most $d^2$ product gates and $d$ sum gates per internal node; hence $W\leq d^2+d+1=O(d^2)$, the theorem's product-region parameter is $I\leq|V|$, and, counting gates and wires, $|C|=|C_P|+|C_Q|=O(|L|d+Id^3)$. The preliminary normalization dynamic program costs $O(|L|qd+Id^3)$ and is subsumed by the claimed bound. Substituting these quantities into Theorem~\ref{thm:pc_fpras} gives the stated running time directly.
\end{proof}

\end{document}
